% =====================================================================
% BRECHAS  ---  SE ESCRIBE ENTRE LOS TRES, AL FINAL
% Sin este apartado el survey es "un catalogo de lecturas" (enunciado).
% Insumo: leer en vertical la columna "Limitacion" de la tabla.
% =====================================================================
\subsection{Identified Research Gaps}
\label{sec:brechas}

\subsubsection*{\textbf{Gaps identified in the literature}} First, horse-racing-specific research is
generally restricted to a single racing venue, jurisdiction, or relatively short
observation period. Davoodi and Khanteymoori~\cite{ejeA_davoodi2010} use one month
of data from a single track, Schumaker~\cite{ejeA_schumaker2018} uses 15 months
from one harness-racing track, and Gupta and Singh~\cite{ejeA_gupta2023,ejeA_gupta2024b}
use three seasons from a single Indian racing operator. Consequently, the
literature provides limited evidence about whether reported models generalize
across venues or racing jurisdictions. Second, Schumaker~\cite{ejeA_schumaker2018}
shows that a race-history feature window developed for greyhound racing performs
poorly when transferred to harness racing, with wagering accuracy decreasing from
50.4\% to 35.4\%. This provides direct evidence that feature configurations should
be validated for the specific racing context rather than transferred without re-
evaluation. Third, Yuan et al.~\cite{ejeB_yuan2025} and Hall and Rasheed~\cite{ejeB_hall2025}
emphasize that comparisons across studies are difficult because differences in
features, data splits, preprocessing, and evaluation procedures can obscure the
effect of the modeling technique itself.

\subsubsection*{\textbf{Patterns inferred by comparison}} First, the horse-racing-specific studies
considered in this project formulate prediction at the individual-horse level,
using a binary or regression target and sorting the resulting predictions
afterward. None of the reviewed horse-racing studies explicitly formulates the
complete race outcome as a learning-to-rank problem. This observation is not
presented as a limitation stated by the authors, but as an inference from the
reviewed corpus; evidence from related racing and ranking domains discussed in the
\ref{sec:ejeB}, suggests that the grouping of competitors within a race may deserve
explicit consideration. Second, only Schumaker~\cite{ejeA_schumaker2018} explicitly
describes a chronological train/test split among the horse-racing studies reviewed.
The remaining studies provide limited information about their temporal splitting
procedures. Therefore, the extent to which their reported results would generalize
to future races is difficult to assess. This observation is treated as a limitation
of the available evidence rather than as proof that the reported experiments contain
temporal leakage.

\subsubsection*{\textbf{What this project addresses}} The present project addresses these
methodological gaps within the limits of the available Hong Kong dataset. We
use a chronological train/validation/test protocol to evaluate performance on
future races and restrict predictive features to information available before
the race, reducing the risk of temporal leakage. We also investigate the race
as the natural grouping unit and evaluate whether a ranking-oriented
formulation is appropriate for the prediction objective, rather than assuming
that independent per-horse classification is sufficient. Feature
configurations are selected and validated using the characteristics of the
Hong Kong dataset rather than being transferred directly from another racing
sub-type. Finally, calibration is considered alongside predictive performance
to distinguish between ranking quality and the reliability of predicted
probabilities

\subsubsection*{\textbf{What remains out of scope}} The present project addresses these
methodological gaps within the limits of the available Hong Kong dataset. We
use a chronological train/validation/test protocol to evaluate performance on
future races and restrict predictive features to information available before
the race, reducing the risk of temporal leakage. We also investigate the race
as the natural grouping unit and evaluate whether a ranking-oriented
formulation is appropriate for the prediction objective, rather than assuming
that independent per-horse classification is sufficient. Feature
configurations are selected and validated using the characteristics of the
Hong Kong dataset rather than being transferred directly from another racing
sub-type. Related studies on physiological and psychological stress
~\cite{ejeA_bartolome2016}, equestrian performance factors~\cite{ejeA_stachurska2019},
and genetic breeding values ~\cite{ejeA_gomez2026} illustrate research directions
that require information we do not have. Finally, calibration is considered
alongside predictive performance to distinguish between ranking quality and the
reliability of predicted probabilities.